% vim: spl=pt
\begin{exercício}{Decomposição do espectro}{ex10}
  Seja \(T \in \bounded(X).\) Mostre que \(\mathbb{C}\) é a união disjunta dos conjuntos \(\rho(T),\) \(\sigma_p(T),\) \(\sigma_c(T),\) e \(\sigma_r(T).\)
\end{exercício}
\begin{proof}[Resolução]
  Notemos que
  \begin{align*}
    \mathbb{C} &= \setc{\lambda \in \mathbb{C}}{\lambda \unity - T \in \mathrm{Inv}(X)} \sqcup \setc{\lambda \in \mathbb{C}}{\lambda \unity - T \notin \mathrm{Inv}(X)}\\
               &= \rho(T) \sqcup \sigma(T),
  \end{align*}
  onde \(\sqcup\) denota a união disjunta. Pelo teorema da aplicação inversa, sabemos que se \(\lambda \unity - T\) for bijetiva, então \(\lambda \unity - T\) tem inversa contínua. Assim,
  \begin{align*}
    \sigma(T) &= \setc{\lambda \in \mathbb{C}}{\ker{(\lambda \unity - T)} = \set{0}\land\ran{(\lambda \unity - T)} \subsetneq X} \sqcup \setc{\lambda \in \mathbb{C}}{\ker{(\lambda \unity - T)} \supsetneq \set{0}}\\
              &= \setc{\lambda \in \mathbb{C}}{\ker{(\lambda \unity - T)} = \set{0} \land \ran{(\lambda \unity - T)} \subsetneq X} \sqcup \sigma_p(T).
  \end{align*}
  Por fim, ou a imagem de \(\lambda \unity - T\) é densa em \(X\) ou não o é, e concluímos que
  \begin{equation*}
      \sigma(T) = \sigma_r(T) \sqcup \sigma_c(T) \sqcup \sigma_p(T),
  \end{equation*}
  logo \(\mathbb{C} = \rho(T) \sqcup \sigma_p(T) \sqcup \sigma_c(T) \sqcup \sigma_r(T).\)
\end{proof}
